%% ============================================================
%% CHAPTER 4 — RQ1: THERMODYNAMIC RESULTS
%% ============================================================
\chapter{Thermodynamic Performance Results (RQ1)}
\label{ch:results_rq1}

This chapter presents all results addressing RQ1. Section~\ref{sec:results_firstlaw}
covers first-law parametric performance; Section~\ref{sec:results_exergy} the
second-law exergy analysis; Section~\ref{sec:results_transitional} the
pseudo-critical heater integration; and Sections~\ref{sec:recup_sensitivity}
and~\ref{sec:plow_sensitivity} the recuperator and condensing-pressure
sensitivity sweeps, together with the warm-climate sensitivity of the sCO$_2$
compressor inlet. The thirteen findings of this chapter differ in several
respects from the Interim Report because of the recuperator correction
described in Section~\ref{sec:method_recup}; the differences are stated
explicitly where they arise.

\section{Model Validation}
\label{sec:results_validation}

SRC validation against NIST Webbook gave a maximum enthalpy deviation of
0.031\,\% (Table~\ref{tab:validation}), well inside the 1\,\% criterion.
CoolProp's CO$_2$ critical point ($T_c = 30.98$\,\textdegree C, $P_c =
7.377$\,MPa) matches NIST REFPROP to within 0.04\,\%. Physical
self-consistency checks at the SRC baseline (400\,\textdegree C TIT,
10\,MPa, 10\,kPa condenser) gave condenser temperature 45.8\,\textdegree C,
BWR 1.31\,\% and $\eta_{th}$ 32.79\,\%, all within expected ranges for
subcritical steam at this condition. The global exergy balance
(Eq.~\ref{eq:exergy_balance}) closes to machine precision for all three
cycles, and the discretised recuperator shows a positive temperature
difference at every node for every case reported (the smallest, 2.7\,K, is
the sCO$_2$ recuperator at the 200\,\textdegree C site, discussed in
Section~\ref{sec:rq2_sitedesign}).

\section{First-Law Parametric Performance}
\label{sec:results_firstlaw}

Figure~\ref{fig:threecycle} shows thermal efficiency and back-work ratio
across 150--600\,\textdegree C turbine inlet temperature for all three cycles
under the sweep rules of Section~\ref{sec:method_thermo}, and
Table~\ref{tab:rq1_results} extracts the values at the three Australian case
study temperatures.

\begin{figure}[H]
\centering
\thesisfig{RQ1_THREE_CYCLE.png}{0.92}{10cm}
\caption{Three-cycle thermal efficiency (top) and back-work ratio (bottom)
versus turbine inlet temperature, 150--600\,\textdegree C. Dashed lines mark
the Australian case study temperatures. Parameters: $\eta_{turbine} = 0.85$,
$\eta_{pump/comp} = 0.80$; SRC: $\le$10\,MPa, 10\,kPa condenser; sCO$_2$:
20/8\,MPa, $\varepsilon = 0.85$, compressor inlet 36\,\textdegree C; tCO$_2$:
$P_{high}(T)$, $P_{low} = 5$\,MPa, $\varepsilon = 0.70$. Recuperators limited
by both streams (Eq.~\ref{eq:recup}).}
\label{fig:threecycle}
\end{figure}

\begin{findingbox}
\textbf{Finding 1 --- Under the modelled assumptions the steam Rankine cycle
has the highest thermal efficiency at every turbine inlet temperature from
150 to 600\,\textdegree C.}
SRC rises from 19.06\,\% to 35.38\,\%; tCO$_2$ from 13.28\,\% to 30.12\,\%;
sCO$_2$ from 5.29\,\% to 30.09\,\%. The two CO$_2$ cycles converge on each
other at 600\,\textdegree C (\keyresult{30.1\,\% for both}) but remain
5.3~percentage points below steam. No crossover occurs within the WHR range.
This reverses the Interim Report's Findings~1 and~2 (sCO$_2$/SRC crossover at
362\,\textdegree C; tCO$_2$ ahead of SRC everywhere), both of which rested on
recuperators that transferred more heat than the hot stream could supply
(Section~\ref{sec:method_recup}). The steam figures are unchanged.
\end{findingbox}

\begin{findingbox}
\textbf{Finding 2 --- sCO$_2$ back-work ratio is prohibitive below
$\approx$300\,\textdegree C.}
At 150\,\textdegree C the compressor absorbs \keyresult{79.8\,\% of turbine
output}; at 200\,\textdegree C, 61.7\,\%; at 300\,\textdegree C, 45.0\,\%;
and still 26.9\,\% at 600\,\textdegree C. The tCO$_2$ pump absorbs 23.1\,\%
at 150\,\textdegree C (higher than the literature-reported 2--5\,\% owing to
$P_{low} = 5$\,MPa), falling to 12.4\,\% at 600\,\textdegree C. The SRC feed
pump stays below 1.5\,\% throughout. BWR, not turbine work, is why the
CO$_2$ cycles trail steam: the gross specific turbine work of the CO$_2$
cycles is comparable to steam per unit of heat input, but a quarter to
three-quarters of it is returned to the compressor.
\end{findingbox}

\begin{findingbox}
\textbf{Finding 3 --- The recuperator gain is bounded by the hot stream.}
With the effectiveness applied to the limiting stream, recuperation at
$\varepsilon = 0.70$ raises the tCO$_2$ efficiency by \keyresult{+3.4, +7.0
and +10.1 percentage points} at the 200, 350 and 500\,\textdegree C site
conditions (12.2$\to$15.7\,\%, 18.0$\to$25.0\,\%, 18.4$\to$28.5\,\%), compared
with the +12.3, +16.8 and +20.3 points obtained at the interim stage with a
cold-side-only effectiveness. The hot-side limit binds because the 5\,MPa
exhaust has a much smaller enthalpy capacity than the near-critical
high-pressure stream it heats.
\end{findingbox}

\begin{table}[H]
\centering
\caption{Thermal efficiency and back-work ratio at the Australian case study
turbine inlet temperatures (parametric sweep rules;
$\varepsilon_{sCO_2} = 0.85$, $\varepsilon_{tCO_2} = 0.70$, $P_{low,tCO_2} = 5$\,MPa).}
\label{tab:rq1_results}
\renewcommand{\arraystretch}{1.35}
\begin{tabularx}{\textwidth}{>{\raggedright\arraybackslash}p{3.8cm}
  >{\centering\arraybackslash}X >{\centering\arraybackslash}X >{\centering\arraybackslash}X}
\toprule
\rowcolor{tableheader}
\textbf{Parameter} &
\textbf{200\,\textdegree C (CS2)} &
\textbf{350\,\textdegree C (CS1)} &
\textbf{500\,\textdegree C (CS3)} \\
\midrule
\rowcolor{rowgray}
SRC $\eta_{th}$ (\%)     & 24.45 & 32.13 & 34.06 \\
sCO$_2$ $\eta_{th}$ (\%) & 12.53 & 22.70 & 27.84 \\
\rowcolor{rowgray}
tCO$_2$ $\eta_{th}$ (\%) & 16.09 & 25.28 & 28.66 \\
\midrule
SRC BWR (\%)              & 0.29  & 1.42  & 1.15 \\
\rowcolor{rowgray}
sCO$_2$ BWR (\%)         & 61.74 & 40.14 & 30.85 \\
tCO$_2$ BWR (\%)         & 19.98 & 18.56 & 14.22 \\
\midrule
\rowcolor{rowgray}
Best cycle ($\eta_{th}$) & SRC & SRC & SRC \\
\bottomrule
\end{tabularx}
\end{table}

\section{Second-Law Analysis}
\label{sec:results_exergy}

\subsection{Basis of Comparison}

Exergy destruction is conventionally tabulated per kilogram of working fluid.
For a comparison of cycles that is misleading, because at a common turbine
inlet temperature the three cycles circulate very different mass flows per
kilowatt: at 400\,\textdegree C the net specific work is 948\,kJ/kg for steam,
108\,kJ/kg for tCO$_2$ and 57\,kJ/kg for sCO$_2$. A kilogram of CO$_2$
destroys far less exergy than a kilogram of steam simply because it does far
less work. Figure~\ref{fig:exergy} therefore shows both bases, and
Table~\ref{tab:exergy} reports destruction per kilojoule of net work, which
is the like-for-like quantity. The Interim Report's statement that sCO$_2$
destroys 95\,\% less exergy than steam compared the per-kilogram totals and
is withdrawn.

\begin{figure}[H]
\centering
\thesisfig{RQ1_EXERGY.png}{1.0}{8cm}
\caption{Component exergy destruction at a common turbine inlet temperature of
400\,\textdegree C (dead state 25\,\textdegree C, 101.325\,kPa), per kilogram
of working fluid (left) and per kilojoule of net work (right). The right-hand
panel is the comparable basis.}
\label{fig:exergy}
\end{figure}

\begin{table}[H]
\centering
\caption{Exergy destruction at 400\,\textdegree C turbine inlet temperature,
in kJ per kJ of net work (per-kilogram totals in the last column).}
\label{tab:exergy}
\renewcommand{\arraystretch}{1.3}
\footnotesize
\begin{tabularx}{\textwidth}{>{\raggedright\arraybackslash}p{1.5cm}
  >{\centering\arraybackslash}X >{\centering\arraybackslash}X >{\centering\arraybackslash}X
  >{\centering\arraybackslash}X >{\centering\arraybackslash}X >{\centering\arraybackslash}p{1.3cm}
  >{\centering\arraybackslash}p{1.3cm} >{\centering\arraybackslash}p{1.7cm}}
\toprule
\rowcolor{tableheader}
\textbf{Cycle} & \textbf{Pump/ comp.} & \textbf{Heater} & \textbf{Turbine} &
\textbf{Cond./ cooler} & \textbf{Recup.} & \textbf{Total} & \textbf{$\eta_{ex}$ (\%)} & \textbf{Total (kJ/kg)} \\
\midrule
\rowcolor{rowgray}
SRC     & 0.002 & 0.396 & 0.167 & 0.134 & ---   & \textbf{0.699} & 58.9 & 663.2 \\
sCO$_2$ & 0.094 & 0.306 & 0.144 & 0.357 & 0.351 & \textbf{1.252} & 44.4 & 71.5 \\
\rowcolor{rowgray}
tCO$_2$ & 0.040 & 0.569 & 0.120 & 0.013 & 0.351 & \textbf{1.092} & 47.8 & 117.7 \\
\bottomrule
\end{tabularx}
\end{table}

\begin{findingbox}
\textbf{Finding 4 --- The steam boiler is the single largest loss in any
cycle, but per unit of work steam remains the most exergy-efficient cycle
at 400\,\textdegree C.}
The boiler destroys \keyresult{375\,kJ/kg, 56.6\,\% of all SRC losses},
because water evaporates at a fixed temperature well below the source. Yet
the SRC destroys only 0.70\,kJ of exergy per kJ of net work against 1.09 for
tCO$_2$ and 1.25 for sCO$_2$, with exergy efficiencies of 58.9, 47.8 and
44.4\,\%. The CO$_2$ cycles' non-isothermal heating does reduce heater
destruction per unit of work (0.31 for sCO$_2$ versus 0.40 for steam), but
the saving is more than consumed by their recuperators (0.35 each) and, for
sCO$_2$, by the compressor and gas cooler (0.09 and 0.36).
\end{findingbox}

\begin{findingbox}
\textbf{Finding 5 --- For the CO$_2$ cycles the recuperator, not the heater,
is the dominant internal irreversibility.}
At 400\,\textdegree C the recuperator accounts for 28\,\% of sCO$_2$
destruction and 32\,\% of tCO$_2$ destruction, as large as the heater in the
sCO$_2$ case. The large temperature differences that a fixed-effectiveness
recuperator tolerates when the two streams have very different $c_p$ are a
genuine second-law cost, and they grow with source temperature
(Section~\ref{sec:recup_sensitivity}).
\end{findingbox}

\subsection{Exergy at the Site Design Points}

Table~\ref{tab:exergy_sites} repeats the analysis at the three site design
points of Section~\ref{sec:rq2_sitedesign}, where the steam cycle operates at
the pinch-limited boiler pressure rather than at 10\,MPa. Two features stand
out. First, with a realistic boiler pressure the tCO$_2$ cycle is the most
exergy-efficient option at 200 and 350\,\textdegree C (42.4 and 47.9\,\%
against 39.0 and 45.4\,\% for steam), and steam leads only at 500\,\textdegree C.
Second, the tCO$_2$ condenser term is \emph{negative} at 200 and
350\,\textdegree C. The condenser rejects heat at 14.3\,\textdegree C, below
the 25\,\textdegree C dead state, so the condensate leaving it carries more
exergy (``cold exergy'') than the vapour entering. The balance closes only
because the model does not charge the cycle for the refrigeration needed to
supply a 14\,\textdegree C sink; that charge is quantified in
Section~\ref{sec:plow_sensitivity}.

\begin{table}[H]
\centering
\caption{Exergy destruction at the site design points (per kJ net work) with
the heat source treated as a reservoir at $T_{src}$.}
\label{tab:exergy_sites}
\renewcommand{\arraystretch}{1.3}
\footnotesize
\begin{tabularx}{\textwidth}{>{\raggedright\arraybackslash}p{1.0cm}>{\raggedright\arraybackslash}p{1.1cm}
  >{\centering\arraybackslash}X >{\centering\arraybackslash}X >{\centering\arraybackslash}X
  >{\centering\arraybackslash}X >{\centering\arraybackslash}X >{\centering\arraybackslash}X
  >{\centering\arraybackslash}p{1.2cm} >{\centering\arraybackslash}p{1.2cm}}
\toprule
\rowcolor{tableheader}
\textbf{Site} & \textbf{Cycle} & \textbf{Pump/ comp.} & \textbf{Heater} & \textbf{Turbine} &
\textbf{Cond./ cooler} & \textbf{Recup.} & \textbf{Total} & \textbf{$\eta_{th}$ (\%)} & \textbf{$\eta_{ex}$ (\%)} \\
\midrule
\rowcolor{rowgray} CS2 & SRC     & 0.000 & 1.010 & 0.165 & 0.387 & --- & 1.562 & 14.4 & 39.0 \\
                   CS2 & sCO$_2$ & 0.298 & 0.807 & 0.390 & 0.724 & 0.016 & 2.235 & 11.4 & 30.9 \\
\rowcolor{rowgray} CS2 & tCO$_2$ & 0.052 & 1.085 & 0.175 & $-$0.140 & 0.187 & 1.360 & 15.7 & 42.4 \\
\midrule
                   CS1 & SRC     & 0.000 & 0.828 & 0.165 & 0.210 & --- & 1.204 & 23.7 & 45.4 \\
\rowcolor{rowgray} CS1 & sCO$_2$ & 0.114 & 0.400 & 0.172 & 0.386 & 0.272 & 1.345 & 22.2 & 42.6 \\
                   CS1 & tCO$_2$ & 0.046 & 0.658 & 0.137 & $-$0.017 & 0.265 & 1.089 & 25.0 & 47.9 \\
\midrule
\rowcolor{rowgray} CS3 & SRC     & 0.001 & 0.681 & 0.166 & 0.146 & --- & 0.994 & 30.8 & 50.2 \\
                   CS3 & sCO$_2$ & 0.075 & 0.261 & 0.117 & 0.335 & 0.440 & 1.228 & 27.6 & 44.9 \\
\rowcolor{rowgray} CS3 & tCO$_2$ & 0.033 & 0.504 & 0.101 & 0.055 & 0.463 & 1.157 & 28.5 & 46.4 \\
\bottomrule
\end{tabularx}
\end{table}

\section{Transitional Effects: Pseudo-Critical Heater Integration}
\label{sec:results_transitional}

Figures~\ref{fig:pinch_200}--\ref{fig:pinch_500} show the tCO$_2$ heater T-h
diagrams at the three site conditions with $\varepsilon = 0$, so that the
pseudo-critical $c_p$ peak lies inside the heater, and
Table~\ref{tab:pinch} summarises the integration.

\begin{figure}[H]
\centering
\thesisfig{tco2_pinch_200C.png}{0.9}{9cm}
\caption{tCO$_2$ heater T-h diagram and $c_p$ profile for CS2 gas processing
(200\,\textdegree C source, 12.8\,MPa, $\varepsilon = 0$). Upper panel: CO$_2$
temperature (gold) and waste-gas temperature (red dashed, linear in heat
duty) versus cumulative heat transferred from the cold end. The approach of
10\,K is reached at \emph{both} terminals; the gap widens to 41.6\,K at a
CO$_2$ temperature of 88\,\textdegree C, just beyond the pseudo-critical
temperature of 57.1\,\textdegree C. Lower panel: $c_p$ along the heater, peak
4.42\,kJ/(kg$\cdot$K).}
\label{fig:pinch_200}
\end{figure}

\begin{figure}[H]
\centering
\thesisfig{tco2_pinch_350C.png}{0.9}{9cm}
\caption{tCO$_2$ heater T-h diagram and $c_p$ profile for CS1 cement
(350\,\textdegree C source, 20.0\,MPa, $\varepsilon = 0$). At 20\,MPa the $c_p$
peak is broader and lower (2.62\,kJ/(kg$\cdot$K) at 76\,\textdegree C); the
pump exit temperature rises to 32.1\,\textdegree C, so the gas cannot be
cooled below 42.1\,\textdegree C. Largest gap 49.0\,K at 145\,\textdegree C.}
\label{fig:pinch_350}
\end{figure}

\begin{figure}[H]
\centering
\thesisfig{tco2_pinch_500C.png}{0.9}{9cm}
\caption{tCO$_2$ heater T-h diagram and $c_p$ profile for CS3 steel
(500\,\textdegree C source, 20.0\,MPa, $\varepsilon = 0$). The large
source-to-sink span drives total heater duty to 701\,kJ/kg; largest gap
64.7\,K at 166\,\textdegree C.}
\label{fig:pinch_500}
\end{figure}

\begin{table}[H]
\centering
\caption{tCO$_2$ heater integration at $\varepsilon = 0$ (10\,K approach,
turbine inlet $T_{src} - 10$\,K) and the corresponding design values at
$\varepsilon = 0.70$.}
\label{tab:pinch}
\renewcommand{\arraystretch}{1.3}
\begin{tabularx}{\textwidth}{>{\raggedright\arraybackslash}p{6.2cm}
  >{\centering\arraybackslash}X >{\centering\arraybackslash}X >{\centering\arraybackslash}X}
\toprule
\rowcolor{tableheader}
\textbf{Parameter} & \textbf{CS2 (200\,\textdegree C)} & \textbf{CS1 (350\,\textdegree C)} & \textbf{CS3 (500\,\textdegree C)} \\
\midrule
\rowcolor{rowgray}
$P_{high}$ (MPa)                                 & 12.8  & 20.0  & 20.0 \\
Pump exit temperature (\textdegree C)            & 24.3  & 32.1  & 32.1 \\
\rowcolor{rowgray}
$T_{pc}$ (\textdegree C) / peak $c_p$ (kJ/kg$\cdot$K) & 57.1 / 4.42 & 76.2 / 2.62 & 75.8 / 2.62 \\
Gas outlet temperature, $\varepsilon = 0$ (\textdegree C) & 34.3 & 42.1 & 42.1 \\
\rowcolor{rowgray}
Heater duty, $\varepsilon = 0$ (kJ/kg)           & 360.5 & 517.3 & 701.3 \\
Minimum $\Delta T$ (K) and location              & 10.0, both ends & 10.0, both ends & 10.0, both ends \\
\rowcolor{rowgray}
Maximum $\Delta T$ (K) at $T_{CO_2}$ (\textdegree C) & 41.6 at 88 & 49.0 at 145 & 64.7 at 166 \\
Effective LMTD, segmented (K)                    & 22.9  & 27.1  & 32.3 \\
\midrule
\rowcolor{rowgray}
Heater inlet, $\varepsilon = 0.70$ (\textdegree C) & 50.2 & 91.1 & 141.1 \\
Gas outlet, $\varepsilon = 0.70$ (\textdegree C)  & 60.2 & 101.1 & 151.1 \\
\rowcolor{rowgray}
Effective LMTD, $\varepsilon = 0.70$ (K)          & 23.2 & 21.3 & 16.1 \\
\bottomrule
\end{tabularx}
\end{table}

\begin{findingbox}
\textbf{Finding 6 --- The tCO$_2$ heater is pinched at both terminals, and the
pseudo-critical $c_p$ peak opens the driving force in the middle rather than
closing it.}
With the gas profile linear in heat duty, the 10\,K approach is reached at
the cold end and at the hot end simultaneously, and the temperature
difference bulges to \keyresult{41.6, 49.0 and 64.7\,K} at CO$_2$ temperatures
of 88, 145 and 166\,\textdegree C, downstream of $T_{pc}$. The mechanism is
that CO$_2$ absorbs a large share of the duty while its temperature is still
low (high $c_p$), so its curve falls below the straight line joining its end
points while the gas, with constant $c_p$, follows that straight line. The
Interim Report's Finding~7 (``the pinch sits mid-exchanger at
129\,\textdegree C'') was an artefact of an index-linear gas profile that
produced a constant 10\,K gap everywhere, and is withdrawn.
\end{findingbox}

\begin{findingbox}
\textbf{Finding 7 --- The heater inlet temperature sets a hard floor on waste
heat recovery, and recuperation raises that floor.}
At $\varepsilon = 0$ the floor is the pump exit: the gas cannot leave below
\keyresult{34.3\,\textdegree C} at 12.8\,MPa or 42.1\,\textdegree C at
20\,MPa. At the design effectiveness of 0.70 the heater inlet rises to 50,
91 and 141\,\textdegree C and the gas can only be cooled to 60, 101 and
151\,\textdegree C. Recuperation therefore trades depth of heat recovery for
cycle efficiency, a trade quantified in Section~\ref{sec:recup_sensitivity}.
\end{findingbox}

\begin{findingbox}
\textbf{Finding 8 --- The pseudo-critical peak sharpens as pressure falls.}
At 12.8\,MPa, $T_{pc} = 57.1$\,\textdegree C with peak $c_p =
4.42$\,kJ/(kg$\cdot$K); at 20\,MPa, $T_{pc} = 76$\,\textdegree C with peak
$c_p = 2.62$\,kJ/(kg$\cdot$K). Lower pressure produces the more non-uniform
profile and therefore the larger departure between segmented and two-point
sizing, but the pump exit temperature is also lower, which deepens heat
recovery. The high-side pressure is thus a trade-off between heater
non-uniformity and recovery depth, in addition to its effect on efficiency.
\end{findingbox}

\section{Recuperator Sensitivity}
\label{sec:recup_sensitivity}

Figure~\ref{fig:recup_sensitivity} and Table~\ref{tab:recup} show the tCO$_2$
cycle at $\varepsilon = 0$, 0.50, 0.70 and 0.85 at each site design point.

\begin{figure}[H]
\centering
\thesisfig{RQ1_RECUPERATOR_SENSITIVITY.png}{1.0}{9cm}
\caption{tCO$_2$ recuperator effectiveness sensitivity at the three site design
points: thermal efficiency, net work per kilogram of waste gas, recuperator
duty and the achievable gas outlet temperature.}
\label{fig:recup_sensitivity}
\end{figure}

\begin{table}[H]
\centering
\caption{tCO$_2$ recuperator sensitivity ($P_{low} = 5$\,MPa, 10\,K approach).
$w_{gas}$ is the net work per kilogram of waste gas (Eq.~\ref{eq:wgas}).}
\label{tab:recup}
\renewcommand{\arraystretch}{1.25}
\footnotesize
\begin{tabularx}{\textwidth}{>{\raggedright\arraybackslash}p{1.0cm}>{\centering\arraybackslash}p{0.9cm}
  >{\centering\arraybackslash}X >{\centering\arraybackslash}X >{\centering\arraybackslash}X
  >{\centering\arraybackslash}X >{\centering\arraybackslash}X >{\centering\arraybackslash}X
  >{\centering\arraybackslash}p{1.6cm}}
\toprule
\rowcolor{tableheader}
\textbf{Site} & $\varepsilon$ & $\eta_{th}$ (\%) & BWR (\%) & $q_{in}$ (kJ/kg) & $Q_{recup}$ (kJ/kg) &
Heater inlet (\textdegree C) & Gas out (\textdegree C) & $w_{gas}$ (kJ/kg gas) \\
\midrule
\rowcolor{rowgray} CS2 & 0    & 12.23 & 20.69 & 360.5 & 0.0   & 24.3  & 34.3  & 22.30 \\
                   CS2 & 0.50 & 14.51 & 20.69 & 304.0 & 56.5  & 44.1  & 54.1  & 23.28 \\
\rowcolor{rowgray} CS2 & 0.70 & 15.67 & 20.69 & 281.4 & 79.1  & 50.2  & 60.2  & 24.10 \\
                   CS2 & 0.85 & 16.68 & 20.69 & 264.5 & 96.0  & 54.2  & 64.2  & 24.91 \\
\midrule
\rowcolor{rowgray} CS1 & 0    & 17.97 & 18.97 & 517.3 & 0.0   & 32.1  & 42.1  & 60.85 \\
                   CS1 & 0.50 & 22.47 & 18.97 & 413.7 & 103.6 & 75.0  & 85.0  & 65.49 \\
\rowcolor{rowgray} CS1 & 0.70 & 24.97 & 18.97 & 372.3 & 145.1 & 91.1  & 101.1 & 68.36 \\
                   CS1 & 0.85 & 27.25 & 18.97 & 341.2 & 176.2 & 104.2 & 114.2 & 70.68 \\
\midrule
\rowcolor{rowgray} CS3 & 0    & 18.39 & 14.44 & 701.3 & 0.0   & 32.1  & 42.1  & 92.60 \\
                   CS3 & 0.50 & 24.62 & 14.44 & 523.7 & 177.6 & 104.8 & 114.8 & 104.33 \\
\rowcolor{rowgray} CS3 & 0.70 & 28.49 & 14.44 & 452.7 & 248.7 & 141.1 & 151.1 & 109.32 \\
                   CS3 & 0.85 & 32.29 & 14.44 & 399.4 & 302.0 & 174.7 & 184.7 & 111.99 \\
\bottomrule
\end{tabularx}
\end{table}

\begin{findingbox}
\textbf{Finding 9 --- Recuperation improves efficiency far more than it
improves recovered work, because every increment of $\varepsilon$ raises the
gas outlet temperature.}
From $\varepsilon = 0$ to 0.85, $\eta_{th}$ rises by 4.5, 9.3 and
13.9~points at CS2, CS1 and CS3, but $w_{gas}$ rises by only
\keyresult{+12\,\%, +16\,\% and +21\,\%}, because the gas outlet temperature
climbs from 34--42\,\textdegree C to 64--185\,\textdegree C. The marginal
benefit also diminishes: the step 0.70$\to$0.85 adds 0.8--2.7\,kJ per
kilogram of gas while roughly doubling the recuperator conductance. BWR is
independent of $\varepsilon$ (20.7, 19.0 and 14.4\,\%), confirming that
recuperation reduces external heat input only and does not alter the pump or
turbine work.
\end{findingbox}

\begin{findingbox}
\textbf{Finding 10 --- At high effectiveness the recuperator approaches the
primary heater in duty, and at low source temperature it approaches the
10\,K approach limit.}
At $\varepsilon = 0.85$ the recuperator transfers 302\,kJ/kg at CS3, 76\,\%
of the external heat input (399\,kJ/kg). At CS2 the same effectiveness leaves
a minimum internal temperature difference of only \keyresult{9.8\,K},
essentially at the design approach, so $\varepsilon = 0.85$ is the practical
ceiling for a 200\,\textdegree C source at 12.8\,MPa. The baseline of 0.70
retains 21--86\,K of internal margin at all three sites.
\end{findingbox}

\section{Low-Side Pressure and Heat-Sink Sensitivity}
\label{sec:plow_sensitivity}

Figure~\ref{fig:plow_sensitivity} and Table~\ref{tab:plow} show the tCO$_2$
cycle at low-side pressures from 3.5 to 7.0\,MPa with $\varepsilon = 0.70$.

\begin{figure}[H]
\centering
\thesisfig{RQ1_PLOW_SENSITIVITY.png}{1.0}{7cm}
\caption{tCO$_2$ low-side pressure sensitivity ($\varepsilon = 0.70$):
condensing temperature, thermal efficiency and back-work ratio. The shaded
band marks the cooling-tower water temperatures typical of the three
Australian sites.}
\label{fig:plow_sensitivity}
\end{figure}

\begin{table}[H]
\centering
\caption{tCO$_2$ low-side pressure sensitivity ($\varepsilon = 0.70$, 10\,K
approach).}
\label{tab:plow}
\renewcommand{\arraystretch}{1.25}
\footnotesize
\begin{tabularx}{\textwidth}{>{\raggedright\arraybackslash}p{1.0cm}>{\centering\arraybackslash}p{1.3cm}
  >{\centering\arraybackslash}X >{\centering\arraybackslash}X >{\centering\arraybackslash}X
  >{\centering\arraybackslash}X >{\centering\arraybackslash}X >{\centering\arraybackslash}p{1.7cm}}
\toprule
\rowcolor{tableheader}
\textbf{Site} & $P_{low}$ (MPa) & $T_{cond}$ (\textdegree C) & $\eta_{th}$ (\%) & BWR (\%) & $w_{net}$ (kJ/kg) & Gas out (\textdegree C) & $w_{gas}$ (kJ/kg gas) \\
\midrule
\rowcolor{rowgray} CS2 & 3.5 & 0.2  & 18.47 & 16.72 & 61.6  & 44.7  & 31.56 \\
                   CS2 & 4.0 & 5.3  & 17.49 & 17.98 & 55.1  & 50.7  & 28.72 \\
\rowcolor{rowgray} CS2 & 5.0 & 14.3 & 15.67 & 20.69 & 44.1  & 60.2  & 24.10 \\
                   CS2 & 6.0 & 22.0 & 13.99 & 23.93 & 34.8  & 67.8  & 20.33 \\
\rowcolor{rowgray} CS2 & 7.0 & 28.7 & 12.36 & 28.88 & 26.3  & 76.5  & 16.80 \\
\midrule
                   CS1 & 3.5 & 0.2  & 27.59 & 15.59 & 117.5 & 80.4  & 81.82 \\
\rowcolor{rowgray} CS1 & 4.0 & 5.3  & 26.67 & 16.68 & 108.5 & 87.6  & 76.96 \\
                   CS1 & 5.0 & 14.3 & 24.97 & 18.97 & 93.0  & 101.1 & 68.36 \\
\rowcolor{rowgray} CS1 & 6.0 & 22.0 & 23.41 & 21.64 & 79.5  & 114.8 & 60.57 \\
                   CS1 & 7.0 & 28.7 & 21.91 & 25.55 & 67.0  & 131.0 & 52.79 \\
\midrule
\rowcolor{rowgray} CS3 & 3.5 & 0.2  & 31.56 & 11.84 & 161.5 & 120.2 & 131.86 \\
                   CS3 & 4.0 & 5.3  & 30.48 & 12.68 & 149.6 & 130.4 & 123.92 \\
\rowcolor{rowgray} CS3 & 5.0 & 14.3 & 28.49 & 14.44 & 129.0 & 151.1 & 109.32 \\
                   CS3 & 6.0 & 22.0 & 26.67 & 16.49 & 111.2 & 172.9 & 95.93 \\
\rowcolor{rowgray} CS3 & 7.0 & 28.7 & 24.95 & 19.49 & 94.9  & 197.7 & 82.97 \\
\bottomrule
\end{tabularx}
\end{table}

\begin{findingbox}
\textbf{Finding 11 --- The thermodynamically attractive condensing pressures
are climatically infeasible in Australia, and the 5\,MPa baseline is itself
marginal.}
$P_{low} = 3.5$\,MPa condenses at 0.2\,\textdegree C and 4.0\,MPa at
5.3\,\textdegree C --- unattainable at any of the three sites. The baseline
5\,MPa condenses at \keyresult{14.3\,\textdegree C}, which requires a sink
near 4\,\textdegree C for a 10\,K approach: colder than cooling-tower water
anywhere in Australia (typically 20--32\,\textdegree C~\cite{BOM2024}) and
achievable only with mechanical chilling. Only $P_{low} \ge 6.5$\,MPa
(condensing above 25\,\textdegree C) is compatible with evaporative cooling
at Gladstone or Sydney outside summer, and no subcritical pressure is
compatible with Pilbara summer conditions.
\end{findingbox}

\begin{findingbox}
\textbf{Finding 12 --- The cost of a cold sink would exceed the benefit of
the cycle at low source temperatures.}
If the 14.3\,\textdegree C condenser at 5\,MPa were served by an electric
chiller, the parasitic load would be $q_{out}/\mathrm{COP}$. With
$q_{out} = 237$, 279 and 324\,kJ/kg and $w_{net} = 44$, 93 and 129\,kJ/kg,
a chiller of COP~4 would consume \keyresult{134\,\%, 75\,\% and 63\,\%} of
the net output at CS2, CS1 and CS3 (90, 50 and 42\,\% at COP~6). Raising
$P_{low}$ to 7\,MPa (28.7\,\textdegree C) removes the need for chilling at
the cost of 3.1--3.5 efficiency points and 23--30\,\% of $w_{gas}$; at CS1
the resulting 52.8\,kJ per kilogram of gas equals that of the pinch-limited
steam cycle (Section~\ref{sec:rq2_sitedesign}).
\end{findingbox}

\begin{findingbox}
\textbf{Finding 13 --- The warm-climate penalty on sCO$_2$ is concentrated at
low source temperatures.}
Raising the compressor inlet from 32 to 40\,\textdegree C (8\,MPa) reduces
sCO$_2$ efficiency from \keyresult{15.7 to 9.2\,\%} at a 200\,\textdegree C
turbine inlet (a 42\,\% relative loss), from 23.6 to 21.3\,\% at
350\,\textdegree C, and from 28.0 to 27.1\,\% at 500\,\textdegree C. The
baseline of 36\,\textdegree C, chosen as 5\,K above the critical temperature
to guarantee single-phase compression with Australian cooling water, gives
12.5, 22.7 and 27.8\,\%. Near-critical compression is the source of the
sCO$_2$ cycle's advantage in the literature and it is exactly what a warm heat
sink takes away.
\end{findingbox}

\section{Comparison with Published Results}
\label{sec:lit_comparison}

Kizilkan~\cite{Kizilkan2020} reported 27.6\,\% for sCO$_2$ and 24.2\,\% for
steam at a 350\,\textdegree C cement source. The present steam value at the
350\,\textdegree C \emph{site design point} (pinch-limited to 0.91\,MPa,
Section~\ref{sec:rq2_sitedesign}) is 23.7\,\%, in close agreement; the
present sCO$_2$ value (22.2\,\% at the site, 22.7\,\% in the sweep) is lower
than Kizilkan's because his configuration is a recompression cycle with a
compressor inlet near 32\,\textdegree C, whereas this study models a simple
recuperated layout with a 36\,\textdegree C inlet and a hot-side-limited
recuperator. Jiang et al.~\cite{Jiang2023} found tCO$_2$ up to 5~points
ahead of sCO$_2$; this study finds tCO$_2$ ahead of sCO$_2$ by 2.6--3.6~points
between 200 and 350\,\textdegree C, narrowing to 0.8~points at
500\,\textdegree C, consistent in direction. Crespi et al.'s~\cite{Crespi2017}
40--50\,\% efficiencies for sCO$_2$ apply above 500\,\textdegree C with
optimised recompression layouts and cold sinks, conditions outside the WHR
range and the Australian climate considered here. The disagreement with the
literature is therefore not with the thermodynamics but with the boundary
conditions, and Findings~11--13 identify which ones.

\section{Summary of RQ1 Results}
\label{sec:rq1_summary}

Under consistent assumptions and a second-law-consistent recuperator, the
steam Rankine cycle has the highest thermal efficiency at every turbine inlet
temperature in the 150--600\,\textdegree C range; tCO$_2$ is second below
500\,\textdegree C and the two CO$_2$ cycles converge above it. Per unit of
net work the SRC is also the most exergy-efficient cycle at a common turbine
inlet temperature, although at realistic pinch-limited boiler pressures the
tCO$_2$ cycle leads at 200 and 350\,\textdegree C. The tCO$_2$ heater is
pinched at both terminals with the driving force widest in the middle, so
the transitional effect enlarges rather than restricts the mean temperature
difference. Recuperation raises efficiency but also the gas outlet
temperature, so it adds far less to the work recovered from a stream than to
the efficiency. Finally, the CO$_2$ cycles' sensitivity to the heat sink is
their defining weakness in Australia: the tCO$_2$ baseline needs a
14\,\textdegree C condenser that would cost more to provide than the cycle
produces at 200\,\textdegree C, and the sCO$_2$ cycle loses two-fifths of its
efficiency at 200\,\textdegree C if its compressor inlet warms from 32 to
40\,\textdegree C.
